\chapter{Accès aux marchés~: plateformes crypto}\label{ch:m3:getting-access-crypto-venues}

Une nouvelle société de trading veut coter sur cinq plateformes crypto. Elle ouvre des comptes,
se connecte et négocie. À la fin de son premier mois, elle constate qu'elle a payé, par dollar échangé,
plusieurs fois ce que paie un teneur de marché établi pour les mêmes transactions, car les frais diminuent avec le volume
et son volume est dilué~; que ses connexions ont été bridées à une fraction du débit dont sa stratégie avait besoin~;
et que la remise sur laquelle elle comptait est allée aux sociétés qui maintenaient leurs cotations affichées
quelques points de pourcentage de temps en plus. L'accès à une plateforme crypto est peu coûteux à obtenir et cher à bien utiliser.
Ce chapitre complète les chapitres d'accès de la série pour les crypto-actifs~: les entités et l'entrée en relation qu'exige une plateforme,
les grilles tarifaires et la manière de se placer au bon palier, les programmes de teneur de marché et ce qu'ils exigent,
la connectivité, ainsi que les dispositifs de crédit et de conservation qui permettent à une société de négocier sans laisser ses actifs sur les plateformes.

\section{Entités, entrée en relation et sous-comptes}

\begin{definition}[Contrôle KYC]\label{def:m3:getting-access-crypto-venues:kyc}
Un \emph{contrôle KYC}\index{contrôle KYC} est la vérification, par une plateforme, de l'identité d'un client, de la propriété
et du contrôle qu'il exerce, de l'origine des fonds et de son statut réglementaire, effectuée avant l'ouverture d'un compte, puis répétée périodiquement et lorsque les faits changent.
\end{definition}

Un compte institutionnel sur une plateforme crypto est ouvert par une entité juridique, et l'entrée en relation de la plateforme demande ce que demanderait celle d'un courtier
(livre~1, chapitre~29)~: les statuts de constitution, l'identifiant d'entité juridique (LEI) de l'entité (livre~2, chapitre~28), les bénéficiaires
effectifs et les dirigeants (la vérification d'entreprise de Kraken demande les documents de constitution, un registre des actions des bénéficiaires effectifs ultimes avec leurs
droits financiers et de vote, un organigramme de propriété et une preuve de l'adresse d'exploitation), l'origine des fonds initiaux, les agréments réglementaires de la société et sa stratégie de trading. La juridiction propre à la plateforme
décide de ceux qu'elle peut accepter~: une plateforme offshore peut exclure les résidents des pays où elle n'est pas agréée, et une
plateforme agréée peut exiger que le client le soit aussi (\cref{ch:m3:the-crypto-trading-business}). Une fois le compte ouvert, le compte principal
est divisé en \omterm{def:m3:centralised-exchanges:subaccount}{sous-comptes} (\cref{ch:m3:centralised-exchanges}) par stratégie, chacun disposant de ses propres \omterm{def:m3:centralised-exchanges:subaccount}{clés d'API} et de ses limites.

\section{Grilles tarifaires, paliers VIP et alignement des paliers}

Les plateformes crypto facturent selon les paliers de volume de Marchés~I, avec des taux maker et taker en points de base, et des paliers fixés sur le
volume de négociation sur trente jours du compte (parfois de l'ensemble de ses \omterm{def:m3:centralised-exchanges:subaccount}{sous-comptes}), et parfois sur les avoirs en
\omterm{def:m3:blockchains-for-traders:key}{jeton} de la plateforme.

\begin{dated}{2026-09}{Deux grilles tarifaires spot complètes}\label{dat:m3:getting-access-crypto-venues:tiers}
La grille spot de Binance va du palier régulier (sous 1~million de dollars de volume sur trente jours, 0.100~\% en maker et en taker) jusqu'au VIP~1
(1~million de dollars, 0.090~\%/0.100~\%), VIP~2 (5~millions de dollars, 0.080~\%/0.100~\%), VIP~3 (20~millions de dollars, 0.040~\%/0.060~\%), VIP~4 (75~millions de dollars,
0.040~\%/0.052~\%), VIP~5 (150~millions de dollars, 0.025~\%/0.031~\%), VIP~6 (400~millions de dollars, 0.020~\%/0.029~\%), VIP~7 (800~millions de dollars, 0.019~\%/0.028~\%) et
VIP~8 (2~milliards de dollars, 0.016~\%/0.025~\%) jusqu'au VIP~9 (4~milliards de dollars, 0.011~\%/0.023~\%), chaque palier mentionnant aussi une détention de BNB. La grille spot de Kraken Pro
va de 0.40~\%/0.80~\% sous 2\,500 dollars de volume sur trente jours, à travers 17 paliers, pour atteindre des frais maker de 0.00~\% à partir de 10~millions de dollars, et
0.00~\%/0.05~\% à partir de 500~millions de dollars.
\end{dated}

Les paliers font dépendre le coût d'une société de la façon dont elle répartit son volume. Les frais baissent avec le volume sur chaque plateforme, de sorte que répartir un volume fixe
sur cinq plateformes laisse chacune sur un palier plus bas que si on le concentrait sur deux.

\begin{omfigure}
\begin{tikzpicture}
\begin{axis}[width=11cm,height=5cm,xmode=log,log basis x=10,
  xlabel={volume sur trente jours, en millions de dollars (échelle logarithmique)},ylabel={frais moyen pondéré, pb},
  tick label style={font=\small},label style={font=\small},
  xmin=0.1,xmax=1000,ymin=0,ymax=26,grid=major,grid style={gray!25},
  xtick={0.1,1,10,100,1000},xticklabels={0.1,1,10,100,1\,000},
  legend columns=2,legend style={font=\footnotesize,at={(0.5,-0.3)},anchor=north,draw=none,fill=none,/tikz/every even column/.append style={column sep=8pt}},
  table/col sep=comma]
\addplot[omThm,very thick,const plot] table[x=vol_m,y=binance_bp]{figdata/markets-3/25-getting-access-crypto-venues/feecurve.csv};\addlegendentry{Binance au comptant}
\addplot[omDef,very thick,dashed,const plot] table[x=vol_m,y=kraken_bp]{figdata/markets-3/25-getting-access-crypto-venues/feecurve.csv};\addlegendentry{Kraken Pro au comptant}
\end{axis}
\end{tikzpicture}

\omcaption{Frais moyens pondérés (70~\% du volume en maker, 30~\% en taker) sur les deux grilles publiées de l'\cref{dat:m3:getting-access-crypto-venues:tiers}, en
fonction du volume sur trente jours. Chaque marche correspond à un palier~: en dessous de 2.5~millions de dollars, les taux de Binance sont les plus bas, au-dessus, ceux de Kraken, dont la commission maker tombe à zéro à partir de
10~millions de dollars. Source des données~: le tutoriel du chapitre.}\label{fig:m3:getting-access-crypto-venues:curve}
\end{omfigure}

\begin{proposition}[La concentration réduit la facture de frais]\label{prop:m3:getting-access-crypto-venues:concentrate}
Si la grille de chaque plateforme est non croissante en fonction du volume et que le volume et la part de maker de la société sont fixés, déplacer tout le volume d'une plateforme vers
une autre qui a déjà un taux plus bas à son volume courant ne fait jamais augmenter la facture totale~; l'économie est au moins égale au volume de l'ancienne plateforme multiplié par
l'écart des taux.
\end{proposition}

\begin{proof}
Le volume déplacé paie désormais le taux de la plateforme d'accueil, au plus son ancien taux là-bas puisque le volume de celle-ci a augmenté~; le volume
déjà présent sur la plateforme d'accueil paie un taux qui n'a pas augmenté~; la facture de l'ancienne plateforme disparaît.
\end{proof}

\begin{omfigure}
\begin{tikzpicture}
\begin{axis}[width=11cm,height=5cm,ybar,bar width=12pt,
  symbolic x coords={Binance,Kraken,C,D,E},xtick=data,
  ylabel={frais, milliers de dollars par mois},
  tick label style={font=\small},label style={font=\small},
  ymin=0,ymax=120,grid=major,grid style={gray!25},
  legend columns=2,legend style={font=\footnotesize,at={(0.5,-0.2)},anchor=north,draw=none,fill=none,/tikz/every even column/.append style={column sep=8pt}},
  table/col sep=comma]
\addplot[fill=omThm!60,draw=omThm,area legend] table[x=venue,y=spread]{figdata/markets-3/25-getting-access-crypto-venues/fees.csv};\addlegendentry{réparti sur cinq plateformes}
\addplot[fill=omDef!60,draw=omDef,area legend] table[x=venue,y=concentrated]{figdata/markets-3/25-getting-access-crypto-venues/fees.csv};\addlegendentry{concentré sur deux plateformes}
\end{axis}
\end{tikzpicture}

\omcaption{Un mois de 600~millions de dollars, dont 70~\% en maker, réparti uniformément sur cinq plateformes (321\,000 dollars de frais) ou concentré sur deux (143\,000 dollars).
Binance et Kraken selon les grilles de l'\cref{dat:m3:getting-access-crypto-venues:tiers}~; plateformes C, D et E selon des grilles illustratives. Source des données~:
le tutoriel du chapitre.}\label{fig:m3:getting-access-crypto-venues:fees}
\end{omfigure}

\begin{definition}[Alignement des paliers]\label{def:m3:getting-access-crypto-venues:match}
Un \emph{alignement des paliers}\index{alignement des paliers} est la proposition d'une plateforme d'accorder à un nouveau client, pour une période d'essai, le palier de frais que ce client détient sur une
plateforme concurrente, sur justificatif de ce palier, afin qu'il n'ait pas à monter d'abord dans la grille.
\end{definition}

L'\omterm{def:m3:getting-access-crypto-venues:match}{alignement des paliers} est la manière dont les plateformes se disputent les flux d'ordres déjà concentrés ailleurs. Pour une société, il modifie l'arithmétique de la proposition~:
une troisième plateforme qui aligne son meilleur palier coûte peu à ajouter, et l'économie tirée de la concentration devient une question de savoir si le palier
aligné tient pendant l'essai.

\begin{dated}{2026-09}{Un alignement de paliers et les programmes de liquidité d'une plateforme}\label{dat:m3:getting-access-crypto-venues:programmes}
Du 15 septembre au 31 octobre 2026, Bitget a offert aux traders VIP d'autres bourses une carte de trente jours pour son palier supérieur sur justificatif d'au moins
50~millions d'USDT de volume sur contrats à terme ou 30~millions de volume spot sur les trente jours précédents. Le programme de teneur de marché publié par Bybit
(septembre 2026) verse des remises spot maker de 0.1 point de base au-delà de 25~millions d'USDT de volume mensuel et de 0.5 point de base à partir d'une part de 0.1~\% de
l'ensemble du volume maker de tous les teneurs de marché~; pondère huit fois les perpétuels sur altcoins dans la part qui fixe les paliers des perpétuels, exige des ordres d'au moins dix
fois la taille minimale, revoit les niveaux chaque mois et peut retirer tout avantage à un teneur de marché qui ne les tient pas~; sa grille spot publique
va de 0.10~\% en maker et en taker à 0.030~\% et 0.045~\% pour 100~millions de volume mensuel. Binance a annoncé le 4 juin 2025 un programme sur altcoins
versant 0.5 et 1 point de base aux fournisseurs de liquidité qui atteignent 0.5~\% et 1~\% du volume maker sur des \omterm{def:m3:blockchains-for-traders:key}{jetons} sélectionnés, à côté d'un programme spot maker versant jusqu'à 0.8.
\end{dated}

\section{Programmes des teneurs de marché et leurs obligations}

Au-delà de la grille publique, les plateformes concluent avec les teneurs de marché des programmes~: remises sur le volume maker, parfois un forfait fixe, en contrepartie des
obligations de cotation du type de celles que le livre~1 décrit pour les teneurs de marché désignés (chapitres~24, 29 et~30), mesurées par la plateforme sur ses propres données.

\begin{definition}[Exigence de disponibilité]\label{def:m3:getting-access-crypto-venues:uptime}
Une \emph{exigence de disponibilité}\index{exigence de disponibilité} est la part minimale des instants échantillonnés, sur une période de mesure, auxquels les cotations
d'un teneur de marché sur un instrument doivent satisfaire les conditions du programme des deux côtés~: une fourchette maximale et une taille minimale dans une bande autour du prix milieu.
\end{definition}

\begin{proposition}[La disponibilité se joue dans les queues de distribution]\label{prop:m3:getting-access-crypto-venues:tails}
Si les cotations d'une stratégie ne respectent pas les conditions du programme chaque fois que la volatilité du marché dépasse un seuil $v^*$, sa disponibilité vaut
$\Pr(v < v^*)$, et relever la disponibilité de $1 - p$ à $1 - p'$ impose de coter jusqu'au quantile de volatilité $1 - p'$~: les heures calmes
ne contribuent en rien à la marge.
\end{proposition}

\begin{proof}
Les cotations satisfont les conditions exactement aux instants où $v < v^*$~; la disponibilité est la fréquence de ces instants, soit un quantile de la distribution de la volatilité.
\end{proof}

\begin{omfigure}
\begin{tikzpicture}
\begin{axis}[width=11cm,height=5cm,
  xlabel={fourchette maximale du programme, points de base},ylabel={disponibilité, \%},
  tick label style={font=\small},label style={font=\small},
  xmin=4,xmax=40,ymin=0,ymax=105,grid=major,grid style={gray!25},
  legend columns=3,legend style={font=\footnotesize,at={(0.5,-0.3)},anchor=north,draw=none,fill=none,/tikz/every even column/.append style={column sep=8pt}},
  table/col sep=comma]
\addplot[omThm,very thick] table[x=max_spread,y=thin20]{figdata/markets-3/25-getting-access-crypto-venues/uptime.csv};\addlegendentry{profondeur réduite de moitié au-delà de 2.0 fois la volatilité normale}
\addplot[omDef,very thick,dashed] table[x=max_spread,y=thin225]{figdata/markets-3/25-getting-access-crypto-venues/uptime.csv};\addlegendentry{au-delà de 2.25 fois}
\draw[omExo,thick,dotted] (axis cs:4,90) -- (axis cs:40,90);
\node[font=\footnotesize,anchor=north west,text=omExo] at (axis cs:20,89) {exigence de 90~\%};
\end{axis}
\end{tikzpicture}

\omcaption{Disponibilité d'une stratégie de cotation face à un programme exigeant 50\,000~USD par côté, en fonction de la fourchette maximale admise. La stratégie
cote à 8 points de base en conditions normales, en élargissant sa fourchette avec la volatilité, réduit de moitié sa taille aux instants volatils et retire ses cotations au-delà de trois fois
la volatilité normale. Maintenir la pleine taille jusqu'à 2.25 fois la volatilité normale au lieu de 2.0 fait passer la disponibilité de 87~\% à 91~\%, au-dessus de l'exigence de 90~\% (en pointillés).
Minutes simulées. Données~: le tutoriel du chapitre.}\label{fig:m3:getting-access-crypto-venues:uptime}
\end{omfigure}

L'exigence s'applique précisément quand coter est dangereux~: une société de trading qui retire ses cotations pendant les minutes volatiles, là où sa
sélection adverse est la pire, perd la remise pour tout le mois. Le programme rémunère la prise de risque dans les queues de distribution, et la société doit décider si la
remise en vaut la peine, instrument par instrument.

\section{Limites de débit et endpoints dédiés}

\begin{definition}[Endpoint dédié]\label{def:m3:getting-access-crypto-venues:endpoint}
Un \emph{endpoint dédié}\index{endpoint dédié} est un point d'entrée réseau vers une plateforme réservé à une catégorie de clients, typiquement les teneurs de marché ou
les comptes des plus hauts paliers, avec des \omterm{def:m3:centralised-exchanges:ratelimit}{limites de débit} plus élevées ou un chemin plus court vers le moteur d'appariement que les endpoints publics.
\end{definition}

Les \omterm{def:m3:centralised-exchanges:ratelimit}{limites de débit} du \cref{ch:m3:centralised-exchanges} fixent le budget de cotation d'une société, et la règle la plus contraignante le détermine. Aux limites
de l'\cref{dat:m3:centralised-exchanges:limits}, la règle des ordres de 100 par 10 secondes autorise dix messages par seconde, mais la règle quotidienne de 200\,000 ne permet
que 2.3 messages par seconde en régime soutenu~: une société qui cote cinq instruments peut modifier chaque cotation environ une fois toutes les deux secondes en moyenne, sauf si elle négocie des limites plus élevées.
Les plateformes publient aussi plusieurs endpoints publics aux caractères différents.

\begin{dated}{2026-09}{Les endpoints publics d'une plateforme}\label{dat:m3:getting-access-crypto-venues:endpoints}
La documentation de l'API spot de Binance recense les endpoints de base api.binance.com, api-gcp.binance.com et api1 à api4.binance.com, en notant que les quatre derniers
«~devraient offrir de meilleures performances mais une stabilité moindre~», un endpoint distinct réservé aux seules données de marché publiques, et une fenêtre de validité des requêtes
(\texttt{recvWindow}) de 5\,000 millisecondes par défaut et de 60\,000 au plus.
\end{dated}

\section{Lignes de prêt, règlement hors plateforme et conservation}

\begin{definition}[Ligne de prêt, collateral mirroring]\label{def:m3:getting-access-crypto-venues:loan}
Une \emph{ligne de prêt}\index{ligne de prêt} est une facilité de crédit accordée par une plateforme ou par sa filiale, qui permet à un client institutionnel d'emprunter des actifs ou de négocier sur marge
au-delà de ses dépôts, en contrepartie d'un collatéral et d'une évaluation de crédit. Le \emph{collateral mirroring}\index{collateral mirroring} est le fait de créditer sur une plateforme un solde de négociation
égal au collatéral que le client conserve chez un dépositaire, qui est le mécanisme du \omterm{def:m3:centralised-exchanges:oes}{règlement hors plateforme}
(\cref{ch:m3:centralised-exchanges}).
\end{definition}

La décision d'accès est donc aussi une décision de crédit dans les deux sens. La société accorde du crédit à la plateforme en y déposant des fonds~; la plateforme accorde du crédit
à la société par des \omterm{def:m3:getting-access-crypto-venues:loan}{lignes de prêt} et des soldes miroirs, et le tarife dans les limites et les décotes qu'elle applique. Une société qui négocie sur cinq plateformes avec
un \omterm{def:m3:centralised-exchanges:oes}{règlement hors plateforme} passant par un seul dépositaire détient l'essentiel de son collatéral au même endroit, compense ses expositions entre plateformes et remplace cinq
expositions aux plateformes par une seule exposition au dépositaire et les soldes non réglés dans l'intervalle.

\section{Tutoriel~: cinq plateformes, un budget}\label{tut:m3:getting-access-crypto-venues:fees}

\begin{tutorial}
\textbf{Objectif.} Étant donnés le volume mensuel et la part maker sur cinq plateformes, ainsi que leurs grilles tarifaires, calculer les frais de chaque plateforme et
l'économie liée à la concentration du volume, la valeur d'un alignement de paliers et la disponibilité qu'une stratégie atteint face aux conditions d'un programme. \textbf{État final~:}
les \cref{fig:m3:getting-access-crypto-venues:fees,fig:m3:getting-access-crypto-venues:uptime} et les chiffres du problème du week-end.
\begin{tutsteps}
\item \textbf{Paliers et frais.} Le palier correspondant à un volume, la facture mensuelle et le palier aligné.
\omcode{code/firm/mmprogram/firm_mmprogram.py}{18}{34}{Paliers de volume, frais mensuels et alignement des paliers.}\label{lst:m3:getting-access-crypto-venues:fees}
\item \textbf{Disponibilité.} Instantanés échantillonnés des cotations propres à la société, contrôlés par rapport à la fourchette et à la profondeur.
\omcode{code/firm/mmprogram/firm_mmprogram.py}{46}{58}{Conformité d'un instantané, disponibilité et remise du programme.}\label{lst:m3:getting-access-crypto-venues:uptime}
\item \textbf{Exécuter} \texttt{spread\_vs\_concentrate()}, \texttt{tier\_match\_value()}, \texttt{programme\_value} et \texttt{fig\_cryptoaccess.py}.
\end{tutsteps}
\textbf{Pour aller plus loin.} Trouver la répartition du volume entre les cinq plateformes qui minimise les frais, sous réserve d'au moins 10~\% sur chacune (pour la résilience)~; ajouter un
forfait fixe du programme~; recalculer la disponibilité lorsque la \omterm{def:m3:centralised-exchanges:ratelimit}{limite de débit}, et non la volatilité, est ce qui rend les cotations périmées.
\end{tutorial}

\section{Construction~: le moniteur des programmes de teneurs de marché}\label{bld:m3:getting-access-crypto-venues:mmprogram}

\begin{build}
\bfield{Rôle} La société miniature doit savoir chaque jour à quel palier elle se trouve sur chaque plateforme et ce que coûte sa répartition, si elle respecte les obligations de chaque
programme, et combien de modifications de cotation ses limites autorisent.
\bfield{Interface} \texttt{VolTier(min\_volume, maker\_bp, taker\_bp)}~; \texttt{tier\_for(volume, tiers)}~; \texttt{monthly\_fees(volume, maker\_share,
tiers, matched)}~; \texttt{Snapshot(mid, bid, ask, bid\_depth, ask\_depth)}~; \texttt{compliant}~; \texttt{uptime}~; \texttt{rebate}~; \texttt{requote\_budget(governor,
instruments)} depuis \texttt{firm.ratelimit}.
\bfield{Règles} Grilles versionnées par date~; volume compté comme la plateforme le compte~; instantanés des cotations propres à la société, échantillonnés comme la plateforme les échantillonne~; une
remise d'un programme n'est comptabilisée que lorsque sa mesure la confirme.
\bfield{Tests d'acceptation} \texttt{code/firm/mmprogram/tests/}~: paliers et facture mensuelle~; un alignement de paliers qui n'aide que s'il est meilleur~; conformité, disponibilité et
la remise tout ou rien~; le budget de recotation fixé par la règle d'ordres la plus contraignante.
\bfield{Pour aller plus loin} Un optimiseur de la répartition du volume entre les plateformes~; des alertes de disponibilité intrajournalières~; la mesure du programme reproduite
à partir des données de marché de la plateforme.
\end{build}

\begin{omsources}
\item Binance, spot fee schedule et spot API documentation (general information, limits)~; Kraken, fee schedule. Consulté en septembre 2026.
\item Copper, page produit de ClearLoop (voir le \cref{ch:m3:centralised-exchanges}).
\item Kraken, support article ``How to get verified for a Business account'' (mis à jour le 24 août 2026)~; Bitget, communiqué de presse du 16 septembre 2026~;
Binance, communiqué de presse du 4 juin 2025 (PR Newswire)~; Bybit, centre d'aide, ``Introduction to the Market Maker Incentive Program'' et
``Bybit Trading Fee Structure'', septembre 2026.
\end{omsources}

\section{Exercices}

\begin{exercise}[$\star$]\label{exo:m3:getting-access-crypto-venues:1}
Sur la grille de Binance, combien coûte un mois de 100 millions de dollars, dont 70~\% en tant que teneur de marché~?
\end{exercise}

\begin{exercise}[$\star$]\label{exo:m3:getting-access-crypto-venues:2}
Que demande le \omterm{def:m3:getting-access-crypto-venues:kyc}{contrôle KYC} d'une plateforme à une société de trading, et pourquoi la plateforme y tient-elle~?
\end{exercise}

\begin{exercise}[$\star$]\label{exo:m3:getting-access-crypto-venues:3}
Pourquoi la règle d'ordres journalière, et non celle de dix secondes, fixe-t-elle le budget soutenu d'une société qui cote aux limites de
l'\cref{dat:m3:centralised-exchanges:limits}~?
\end{exercise}

\begin{exercise}[$\star\star$]\label{exo:m3:getting-access-crypto-venues:4}
La plateforme C s'aligne sur le palier que la société détient chez Binance, pour un volume de 120 millions de dollars. Que rapporte cet alignement sur un mois~?
\end{exercise}

\begin{exercise}[$\star\star$]\label{exo:m3:getting-access-crypto-venues:5}
Un programme verse 0.5 point de base sur le volume en tant que teneur de marché, avec une disponibilité de 90~\%. Sur 210 millions de dollars de volume de teneur de marché, que vaut-il, et combien coûte
le fait de manquer l'exigence d'un point~?
\end{exercise}

\begin{exercise}[$\star\star$]\label{exo:m3:getting-access-crypto-venues:6}
Pourquoi une société garderait-elle 10~\% de son volume sur chacune de cinq plateformes, alors que le concentrer coûte moins cher~?
\end{exercise}

\begin{exercise}[$\star\star\star$]\label{exo:m3:getting-access-crypto-venues:7}
\emph{Programmation.} Avec \texttt{programme\_value}, calculer la disponibilité pour des écarts maximaux de 10 et de 20 points de base lorsque la profondeur est divisée par deux au-delà de 2.5 fois la
volatilité normale.
\end{exercise}

\begin{exercise}[$\star\star\star$]\label{exo:m3:getting-access-crypto-venues:8}
\emph{Repérer l'erreur.} «~Nous atteindrons le palier VIP plus vite en négociant avec nous-mêmes sur deux \omterm{def:m3:centralised-exchanges:subaccount}{sous-comptes}.~»
\end{exercise}

\section{Problème~: cinq plateformes, un seul budget}

\begin{problem}[{Problème du week-end --- répartition d'un mois de volume}]\label{pb:m3:getting-access-crypto-venues:1}
Une société échange 600 millions de dollars par mois, dont 70~\% en tant que teneur de marché. Elle peut répartir le volume à parts égales entre Binance, Kraken et trois autres plateformes (C, D et E, sur les
grilles illustratives du tutoriel), ou le concentrer sur Binance et Kraken. Une plateforme propose un programme qui verse 0.5 point de base sur le volume de teneur de marché, pour une
disponibilité de 90~\% avec un écart maximal de 20 points de base et 50\,000 dollars de chaque côté.

\medskip\noindent\textbf{Partie I --- La facture.}
\begin{enumerate}
\item Que coûte chaque plateforme lorsque le volume est réparti à parts égales, et au total~?
\item Que coûtent les deux plateformes lorsque le volume est concentré~?
\item Quelle est l'économie mensuelle~?
\item Quels paliers des plateformes changent, et pourquoi Kraken est-elle si bon marché pour un teneur de marché~?
\item Que changerait un alignement de palier sur la plateforme C~?
\end{enumerate}

\medskip\noindent\textbf{Partie II --- Le programme.}
\begin{enumerate}[resume]
\item Sur 300 millions de dollars à la plateforme du programme, que vaut la remise~?
\item Quelle disponibilité la stratégie de la société atteint-elle, et se qualifie-t-elle~?
\item Quelle modification de la stratégie la fait passer au-delà de 90~\%~?
\item Que coûte cette modification en risque~?
\item Comment l'exigence interagit-elle avec les \omterm{def:m3:centralised-exchanges:ratelimit}{limites de débit} de la plateforme~?
\end{enumerate}

\medskip\noindent\textbf{Partie III --- Au-delà des frais.}
\begin{enumerate}[resume]
\item Quelle exposition aux plateformes chaque répartition crée-t-elle~?
\item Comment le \omterm{def:m3:centralised-exchanges:oes}{règlement hors plateforme} modifierait-il la comparaison~?
\item Que perd la société en renonçant à trois plateformes~?
\item Comment doit-elle mesurer où ses transactions sont exécutées au mieux~?
\item À quelle fréquence doit-elle revoir la répartition~?
\end{enumerate}

\medskip\noindent\textbf{Partie IV --- Jugement.}
\begin{enumerate}[resume]
\item Concentrer le volume est-il bon pour le marché~?
\item Pourquoi les plateformes versent-elles des remises plutôt que de baisser les frais pour tous~?
\item Quel est le nombre de plateformes approprié pour une nouvelle société~?
\item Énoncer le \emph{résultat nommé}~: l'économie mensuelle de la concentration sur deux plateformes, et la disponibilité que la société doit maintenir, et maintient actuellement, pour la remise.
\item En une phrase~: que vend une plateforme de crypto-actifs à un teneur de marché~?
\end{enumerate}
\end{problem}

\section{Questions d'entretien}

\begin{interviewq}[$\star$ \iqroles{trader}]\label{iq:m3:getting-access-crypto-venues:1}
Comment les paliers de volume changent-ils l'endroit où vous devez négocier~?
\end{interviewq}

\begin{interviewq}[$\star$ \iqroles{risque}]\label{iq:m3:getting-access-crypto-venues:2}
Que vérifieriez-vous avant de déposer 20 millions de dollars sur une nouvelle plateforme~?
\end{interviewq}

\begin{interviewq}[$\star\star$ \iqroles{développeur}]\label{iq:m3:getting-access-crypto-venues:3}
Comment concevez-vous un système de cotation qui respecte une \omterm{def:m3:getting-access-crypto-venues:uptime}{exigence de disponibilité} sans dépasser les \omterm{def:m3:centralised-exchanges:ratelimit}{limites de débit}~?
\end{interviewq}

\begin{interviewq}[$\star\star$ \iqroles{trader}]\label{iq:m3:getting-access-crypto-venues:4}
Une remise de teneur de marché vaut-elle la peine d'être prise si elle exige de coter pendant les périodes de forte volatilité~?
\end{interviewq}

\begin{interviewq}[$\star\star$ \iqroles{chercheur}]\label{iq:m3:getting-access-crypto-venues:5}
Comment mesureriez-vous si le volume d'une plateforme est réel avant d'adhérer à son programme~?
\end{interviewq}

\begin{interviewq}[$\star\star\star$ \iqroles{développeur, trader}]\label{iq:m3:getting-access-crypto-venues:6}
Construisez le rapport quotidien dont un responsable du trading a besoin sur les cinq relations de la société avec ses plateformes.
\end{interviewq}
